% !TeX root = ../navier-stokes-manuscript.tex
% ============================================================
% 2.1 Vortex Stretching
% ============================================================

Axial stretching lengthens a narrow material vortex-tube segment while its volume stays fixed, so its cross-section contracts, its rotating fluid moves toward the axis, and the interior spins faster.
The velocity field carrying this segment obeys the momentum equation
\[
  \underbrace{\partial_t u}_{\text{local acceleration}}
  +
  \underbrace{(u\cdot\nabla)u}_{\text{nonlinear transport}}
  =
  \underbrace{-\nabla p}_{\text{pressure}}
  +
  \underbrace{\nu\Delta u}_{\text{viscous diffusion}}.
\]
Taking its curl removes pressure and turns nonlinear transport into material advection of vorticity, with strain and viscosity changing the rotation carried by the segment.

\subsection{Vortex Stretching}\label{sub:ThePerfectlyStretchingVortex}

The segment is locally axisymmetric.
Its strain is locally uniform across the material segment and aligned with its axis.
If \(e\) is that axis and \(L(t)\) is the segment's length, then
\[
  S:=\frac12\bigl(\nabla u+(\nabla u)^{\mathsf T}\bigr),
  \qquad
  Se=\sigma(t)e,
  \qquad
  \sigma(t)>0,
  \qquad
  \frac{\dot L(t)}{L(t)}
  =
  e\cdot Se
  =
  \sigma(t).
\]
Incompressibility keeps the material volume \(\mathcal V\) fixed as the segment lengthens.
With effective transverse area \(A(t):=\mathcal V/L(t)\) and effective radius \(r_{\mathrm{eff}}(t):=\sqrt{A(t)/\pi}\),
\[
  \nabla\cdot u=0,
  \qquad
  \frac{d}{dt}(A L)=0,
  \qquad
  \frac{\dot A}{A}=-\sigma(t),
  \qquad
  \frac{\dot r_{\mathrm{eff}}}{r_{\mathrm{eff}}}
  =
  -\frac{\sigma(t)}{2}.
\]
The shrinking radius measures the contraction as the material deformation moves rotating fluid inward.
With \(\omega:=\nabla\times u\), the vorticity carried along the axis satisfies
\[
  \frac{D\omega}{Dt}
  =
  S\omega+\nu\Delta\omega,
  \qquad
  \frac{D}{Dt}
  :=
  \partial_t+u\cdot\nabla,
\]
and under the alignment \(\omega=|\omega|e\),
\[
  S\omega
  =
  \sigma(t)\omega,
  \qquad
  \frac12\frac{D|\omega|^2}{Dt}
  =
  \sigma(t)|\omega|^2
  +
  \nu\,\omega\cdot\Delta\omega.
\]
Viscosity may reinforce or oppose the stretching-driven spin-up, and squared vorticity grows only where the stretching and viscous terms together are positive.
The energy identity proved in \Cref{prop:ClassicalKineticEnergyIdentity} controls velocity in \(L^2\), while vorticity occupies the antisymmetric part of the velocity gradient:
\[
  \norm{u(t)}_2
  \leq
  \norm{u_0}_2
  <\infty,
  \qquad
  \left|\frac12\bigl(\nabla u-(\nabla u)^{\mathsf T}\bigr)\right|_{\mathrm F}
  =
  \frac{|\omega|}{\sqrt2}
  \leq
  |\nabla u|_{\mathrm F}.
\]
Finite kinetic energy does not prevent that rotational gradient from concentrating inside the narrowing segment.
At finer width, its vorticity remains attached to velocity, pressure, transport, and viscosity; the whole field changes scale together.

\divider{\NavierStokes{} Scaling}

At a time \(t_0<T_*\), write \(\ell:=r_{\mathrm{eff}}(t_0)\).
For \(\lambda>1\), define the rescaled velocity and pressure by
\[
  u_\lambda(x,t)
  :=
  \lambda u(\lambda x,\lambda^2t),
  \qquad
  p_\lambda(x,t)
  :=
  \lambda^2p(\lambda x,\lambda^2t).
\]
On the rescaled interval, \(u_\lambda\) is another \NavierStokes{} solution, and at time \(\lambda^{-2}t_0\) its corresponding segment has effective radius \(\lambda^{-1}\ell\); it is not a later state of the opening material segment.
Every momentum term changes with the field:
\[
\begin{aligned}
  \partial_tu_\lambda(x,t)
  &=
  \lambda^3(\partial_tu)(\lambda x,\lambda^2t),
  \\
  (u_\lambda\cdot\nabla)u_\lambda(x,t)
  &=
  \lambda^3\bigl((u\cdot\nabla)u\bigr)(\lambda x,\lambda^2t),
  \\
  \nabla p_\lambda(x,t)
  &=
  \lambda^3(\nabla p)(\lambda x,\lambda^2t),
  \\
  \nu\Delta u_\lambda(x,t)
  &=
  \lambda^3\nu(\Delta u)(\lambda x,\lambda^2t).
\end{aligned}
\]
Nonlinear transport and viscosity acquire the same cubic factor, so neither gains relative ground.
The same scale change separates the Sobolev norms by their order.

\divider{Critical Height}

Let \(\Lambda:=(-\Delta)^{1/2}\).
Since
\[
  \widehat{u_\lambda}(\xi,t)
  =
  \lambda^{-2}\widehat u(\lambda^{-1}\xi,\lambda^2t),
\]
a change of variables gives
\[
  \norm{u_\lambda(t)}_{\dot H^s}^2
  =
  \lambda^{2s-1}
  \norm{u(\lambda^2t)}_{\dot H^s}^2.
\]
At \(s=0\), kinetic energy loses one power of \(\lambda\).
At \(s=1\), the first derivatives gain one.
At \(s=\tfrac12\), the exponent vanishes, and the scale-neutral half-derivative quantity \(H\), the critical height, is
\[
  H(t)
  :=
  \frac12\norm{u(t)}_{\dot H^{1/2}}^2
  =
  \frac12\norm{\Lambda^{1/2}u(t)}_2^2.
\]
Let \(H_\lambda\) denote the same quantity for \(u_\lambda\).
Then
\[
  \norm{u_\lambda(t)}_2^2
  =
  \lambda^{-1}\norm{u(\lambda^2t)}_2^2,
  \qquad
  H_\lambda(t)=H(\lambda^2t),
  \qquad
  \norm{\nabla u_\lambda(t)}_2^2
  =
  \lambda\norm{\nabla u(\lambda^2t)}_2^2.
\]
H remains fixed under the spatial scaling, but that scale neutrality leaves its direction in time unresolved.

\divider{Nonlinear Production versus Viscous Dissipation}

Nonlinear transport changes \(H\) through the signed production
\[
  \mathcal P_{1/2}[u](t)
  :=
  -\operatorname{Re}
  \pair
    {\Lambda^{1/2}u(t)}
    {\Lambda^{1/2}\bigl((u(t)\cdot\nabla)u(t)\bigr)}_{L^2}.
\]
Incompressibility removes pressure, while viscosity subtracts the nonnegative dissipation \(\nu\norm{\Lambda^{3/2}u(t)}_2^2\).
Thus
\[
  H'(t)
  +
  \nu\norm{\Lambda^{3/2}u(t)}_2^2
  =
  \mathcal P_{1/2}[u](t).
\]
H rises when production exceeds dissipation and falls when dissipation exceeds production.
Under the same rescaling,
\[
  \mathcal P_{1/2}[u_\lambda](t)
  =
  \lambda^2\mathcal P_{1/2}[u](\lambda^2t),
  \qquad
  \nu\norm{\Lambda^{3/2}u_\lambda(t)}_2^2
  =
  \lambda^2\nu\norm{\Lambda^{3/2}u(\lambda^2t)}_2^2.
\]
Both actions scale quadratically, leaving unknown the time needed for the nonlinear field to form a smaller structure.
Viscosity has a separate linear heat clock.

\divider{The Heat Clock}

For the heat equation \(v_t=\nu\Delta v\),
\[
  \widehat v(\xi,t+\delta t)
  =
  e^{-\nu|\xi|^2\delta t}\widehat v(\xi,t).
\]
A scale-localized velocity structure of width \(r\) occupies Fourier frequencies \(|\xi|\simeq r^{-1}\).
Its linear viscous change becomes order one when
\[
  \nu|\xi|^2\delta t\simeq1,
\]
so its heat comparison time is
\[
  \tau_\nu(r)
  :=
  \frac{r^2}{\nu}.
\]
This linear comparison time does not determine the nonlinear lifetime.
Halving the width quarters the heat time.
More generally,
\[
  \tau_\nu(\lambda^{-1}r)
  =
  \lambda^{-2}\tau_\nu(r).
\]
The same width-time scaling appears in the full \NavierStokes{} rescaling, though the nonlinear formation time remains unknown.

\Needspace{18\baselineskip}
\divider{The Zeno Cascade}

The dyadic widths and their heat times are
\[
  r_n:=2^{-n}r_0,
  \qquad
  \tau_n:=\frac{r_n^2}{\nu},
\]
and they satisfy
\[
  \tau_n
  =
  \frac{r_0^2}{\nu}4^{-n},
  \qquad
  \sum_{n=0}^{\infty}\tau_n
  =
  \frac{4r_0^2}{3\nu}
  <
  \infty.
\]
This finite arithmetic permits a schedule, but it does not establish one field realizing it.
If one \NavierStokes{} field exhibits localized structures at widths
\[
  r_0>r_1>r_2>\cdots\longrightarrow0
\]
and at sampled times
\[
  t_0<t_1<t_2<\cdots\uparrow T_*<\infty,
  \qquad
  t_{n+1}-t_n\asymp\tau_n,
\]
with comparison constants independent of \(n\), then its possible history has remaining time
\[
  T_*-t_n
  \asymp
  \sum_{k=n}^{\infty}\tau_k
  =
  \frac{4r_n^2}{3\nu}.
\]
The same field would pass through smaller localized structures while both its sampled gaps and its remaining time collapse on the scale \(r_n^2/\nu\).
Even along that possible history, the first and successive Eulerian time derivatives of the same field at one fixed spatial point remain uncontrolled.
